% Generated draft: requires mathematical and source review.
\documentclass[11pt]{article}
\usepackage[margin=25mm]{geometry}
\usepackage{fontspec}
\setmainfont{Latin Modern Roman}
\usepackage{amsmath,amssymb,mathtools,unicode-math}
\setmathfont{Latin Modern Math}
\usepackage{enumitem,needspace,xurl,hyperref}
\setlength{\parindent}{0pt}
\setlength{\parskip}{5pt}
\setlength{\emergencystretch}{3em}
\newcommand{\N}{\mathbb{N}}
\newcommand{\Z}{\mathbb{Z}}
\newcommand{\Q}{\mathbb{Q}}
\newcommand{\R}{\mathbb{R}}
\newcommand{\C}{\mathbb{C}}
\newcommand{\F}{\mathbb{F}}
\newcommand{\A}{\mathbb{A}}
\newcommand{\PP}{\mathbb{P}}
\newcommand{\E}{\mathbb{E}}
\newcommand{\eps}{\varepsilon}
\newcommand{\dd}{\,\mathrm{d}}
\begin{document}
\section*{UnsolvedMath ID: 800078}
\textbf{Source problem code:} catalogue entry
\subsection{Definitions and mathematical statement}
A group \(G\) is \emph{simple} if it is nontrivial and its only normal subgroups are \(\{1\}\) and \(G\). A \emph{left-invariant total order} on \(G\) is a total order \(<\) satisfying
\[
\forall g,a,b\in G,\qquad a<b\Longrightarrow ga<gb.
\]
A group admitting such an order is left-orderable. A total order is \emph{bi-invariant} if it is invariant under both left and right multiplication.

A \(K(G,1)\) is a connected CW complex \(X\) with \(\pi_1(X)\cong G\) and \(\pi_k(X)=0\) for every \(k\geq2\). For a nonnegative integer \(m\), the group \(G\) has \emph{type \(F_m\)} if some \(K(G,1)\) has a finite \(m\)-skeleton. It has type \(F_\infty\) if it has type \(F_m\) for every nonnegative integer \(m\).

\textbf{Question (conditional on the meaning of orderable).} Interpreting the source's word ``orderable'' as ``left-orderable'', is the following assertion true?
\[
\forall n\in\Z,\quad n\geq2\Longrightarrow
\exists G\quad
\left\{
\begin{array}{l}
G\text{ is simple and left-orderable},\\
G\text{ has type }F_n,\\
G\text{ does not have type }F_{n+1}.
\end{array}
\right.
\]
Here \(G\) may depend on \(n\). The last condition means that no \(K(G,1)\) has a finite \((n+1)\)-skeleton; it is not merely a condition on the particular model witnessing type \(F_n\).

The supplied record does not specify whether ``orderable'' means left-orderable or bi-orderable. If bi-orderability is intended, replace ``left-orderable'' in the displayed assertion by ``bi-orderable''. The supplied evidence does not resolve this qualification.
\subsection{Short English statement}
For every integer \(n\geq2\), does there exist a simple orderable group with a classifying space having a finite \(n\)-skeleton, but with no classifying space having a finite \((n+1)\)-skeleton? The intended invariance convention for the order requires confirmation.
\subsection{Sources}
\begin{itemize}
\item Local UnsolvedMath registry, permanent ID \(800078\), exact source problem code \texttt{catalogue entry}, titled ``Foliations and Diffeomorphism Groups --- Question 19''. The record attributes the problem to Nariman and Yazdi (2025) and labels it open; that label is supplied registry information, not an independently verified status assessment.
\item Nariman and Yazdi, \emph{Problem List on Foliations and Diffeomorphism Groups} (2025), Question 19. Supplied source URL: \url{https://arxiv.org/abs/2503.18274}. The provenance records access on 2026-10-16 and extraction method \texttt{arxiv\_latex}; these are supplied provenance details, not a claim of independent consultation.
\item The supplied problem text cites the reference labels \texttt{hyde2019finitely}, \texttt{hyde2023finitely}, \texttt{hyde2023two}, and \texttt{bon2020groups} for constructions of simple orderable groups. Full bibliographic entries for these labels are not supplied.
\item As historical motivation, the supplied 2025 problem text identifies the Hyde--Lodha examples cited as \texttt{hyde2023finitely} as the only finitely presented examples among the constructions it discusses, and reports that those examples have type \(F_\infty\). These are source-reported qualifications, not a new literature assessment.
\end{itemize}
\end{document}
